\subsection{Pipeline implementation}
\label{sec:impl}

\quantmsdiann{} is implemented in Nextflow DSL2 ($\geq$25.10.4) and follows nf-core community standards~\cite{ewels2020nfcore} for pipeline structure, testing, documentation, and continuous integration.
Each computational step is packaged as a versioned container image (Docker, Singularity, or Apptainer); upstream community tools (\texttt{DIA-NN}~1.8.1, ThermoRawFileParser~\cite{hulstaert2020}, pridepy~\cite{kamatchinathan2025pridepy}, qpx) are pulled from BioContainers, while pipeline-specific images (newer \texttt{DIA-NN} releases and \texttt{wiffconverter}) are hosted on \texttt{ghcr.io/bigbio}.
The repository at \url{https://github.com/bigbio/quantmsdiann} is organised as a standard nf-core pipeline with main workflow (\texttt{main.nf}), subworkflows, modules, test data and CI configurations, and documentation.
Pipeline version v2.1.0 (release codename ``Sao Pablo'', 2026-05-08) was used for the analyses reported in this paper.

\subsection{SDRF parsing and configuration generation}
\quantmsdiann{} accepts SDRF~\cite{dai2021sdrf} as primary input.
SDRF parsing is performed using the sdrf-pipelines Python package developed within the quantms project, which validates the SDRF against the proteomics-specific subset of EFO, PSI-MS, and PRIDE ontology terms and emits a normalised per-sample channel for downstream Nextflow processes.
\texttt{DIA-NN} configuration generation extracts enzyme, fixed and variable modification, instrument, precursor mass tolerance, fragment mass tolerance, and missed-cleavage settings from the SDRF and writes a \texttt{DIA-NN} command-line configuration.
Where SDRF columns are absent, \texttt{DIA-NN} defaults appropriate for the inferred instrument are applied.

\subsection{Raw file handling and \texttt{DIA-NN} execution}
Raw inputs are dispatched to format-specific handlers in the file-preparation subworkflow, each conversion or decompression running as an independent, parallelised Nextflow process.
Thermo \texttt{.raw} files are read natively by \texttt{DIA-NN}~$\geq$~2.1.0 (which added Linux native-reader support) or converted to mzML with ThermoRawFileParser~\cite{hulstaert2020} otherwise (\texttt{--mzml\_convert}); Bruker \texttt{.d} folders are ingested natively, with compressed forms (\texttt{.d.tar}/\texttt{.zip}/\texttt{.gz}) decompressed first; and Sciex \texttt{.wiff}~/~\texttt{.wiff2} are converted to mzML with the open-source \texttt{wiffconverter} (\url{https://github.com/bigbio/wiffconverter}), removing the historical dependency on a proprietary ProteoWizard installation.
Generic \texttt{.dia} and pre-converted \texttt{.mzML} inputs pass through unchanged, with optional OpenMS re-indexing (\texttt{--reindex\_mzml}).
\texttt{DIA-NN}~\cite{demichev2020} is then invoked at the version chosen via \texttt{--diann\_version}: the ProteoBench analyses use \diannversion{1.8.1}, \diannversion{2.1.0}, \diannversion{2.2.0}, \diannversion{2.3.2}, and \diannversion{2.5.0} from the same SDRF, and the cell-line reanalyses use \diannversion{2.5.0} unless otherwise stated.

\subsection{Spectral-library generation}
\quantmsdiann{} runs \texttt{DIA-NN} library-free, generating the spectral library in silico from the input FASTA rather than supplying an external experimental library: \texttt{DIA-NN}'s deep-learning predictor builds a predicted library from the protein sequences, each run is searched against it, a cohort-wide empirical library is assembled from the confident first-pass identifications, and a final pass re-searches and quantifies against that refined library.
The source FASTA was the UniProt human reference proteome~\cite{uniprot2025} (UP000005640, \placeholder{release}), with match-between-runs enabled at the cohort level.
In ProteoBench's submission taxonomy this is the predicted-library strategy (\texttt{predictors\_library}~$=$~DIANN); the ProteoBench comparison (Section~\ref{sec:proteobench}) is therefore restricted to community submissions that used the same predicted-from-FASTA strategy.

\subsection{Quality control and \qpx{} export}
Per-run and per-cohort quality control reports are generated by \texttt{pmultiqc}~\cite{larrea2025pmultiqc}, a proteomics extension of MultiQC~\cite{ewels2016multiqc} that summarises \texttt{DIA-NN}-specific metrics including peptide and protein counts, missed cleavages, charge-state distributions, precursor and fragment mass accuracy, and per-run identification yield.
The \qpx{} archive is produced by the \texttt{qpxc} CLI from the \texttt{qpx} package (\url{https://github.com/bigbio/qpx}), which writes Parquet files for identifications (psm, peptide, protein-group, feature), metadata (sample, run, ontology), and provenance, and optionally exports MuData (\texttt{.h5mu}) containers for scverse-based~\cite{virshup2023scverse} downstream analysis.

\subsection{Datasets benchmarked}
\label{sec:datasets}
The benchmark and reanalysis datasets used in this study (ProteoBench DIA modules~5/7/9/10, the three per-cohort cell-line reanalyses PXD003539, PXD004701, PXD030304, the two additional pan-cohort datasets PXD041421 and PXD017199, and the PXD071075 scalability sweep) are summarised together with their instrument, acquisition mode, and deposition year in Additional file~1: Table~S1.
For each dataset, an SDRF was prepared from the original sample metadata and is provided as part of the Additional file~1 deposition.
For each cell-line reanalysis cohort, the original published quantification was downloaded from the PRIDE deposition and used as the comparator, thresholded at 1\% precursor-level FDR for both \quantmsdiann{} and the original analyses.
The conservative contaminant/decoy filter and the like-for-like-threshold rule applied to the original-paper headline counts are described in Additional file~1.

\subsection{Compute environment and benchmarking}
Runtime benchmarks were performed on the EMBL-EBI HPC cluster (LSF scheduler), with per-task CPU and memory allocations dispatched by Nextflow.
Memory profiling was performed using the Nextflow trace report (\texttt{-with-trace}).

